\section{All dimensions and finite-volume sets}\label{sec:dimensions}

Theorem~\ref{thm:normalized-comparison} already gives the comparison for
bounded sets of finite ambient perimeter in dimensions at least four.
We first obtain the same statement in dimensions two and three by smooth
approximation. A single cutoff argument then removes boundedness, and
metric rescaling restores every negative curvature bound.

\subsection{The classical lower-dimensional inequalities}

We first check that the classical smooth comparisons allow disconnected
domains. For $b\ge0$ and $r=\mathcal V_b^{-1}(V)>0$,
\[
 \mathcal I_b'(V)=m\frac{\cs_b(r)}{\sn_b(r)}>0,
 \qquad
 \frac{d}{dr}\frac{\cs_b(r)}{\sn_b(r)}=-\frac1{\sn_b(r)^2}<0.
\]
Thus $\mathcal I_b$ is concave. Since $\mathcal I_b(0)=0$, concavity
between $0$ and $u+v$ gives
\[
 \mathcal I_b(u)\ge\frac{u}{u+v}\mathcal I_b(u+v),\qquad
 \mathcal I_b(v)\ge\frac{v}{u+v}\mathcal I_b(u+v).
\]
Adding proves subadditivity. Consequently, comparisons for the finitely
many connected components of a relatively compact smooth domain imply
the comparison at its total volume.

In dimension three, Kleiner's Theorem~2~\cite[p.~38]{Kleiner1992} gives
the model comparison for every upper sectional-curvature bound
$\kappa\le0$ on a complete simply connected manifold, for arbitrary
compact smooth domains. Apply this result to each connected component
and use subadditivity if necessary.

In dimension two, the classical Weil--Bol inequality
\cite{Weil1926,Bol1941}, in the smooth disk formulation of Chavel
and Feldman \cite[Isoperimetric Inequality~(I), p.~83]{ChavelFeldman1980},
gives, for a smooth disk of area $V$ and boundary length $L$ under
$K\le\kappa\le0$,
\[
 L^2\ge4\pi V-\kappa V^2;
\]
This disk statement suffices for arbitrary relatively compact smooth
domains. The Cartan--Hadamard theorem identifies the surface with $\R^2$.
Direct evaluation of the two-dimensional model functions gives
$\mathcal I_{\sqrt{-\kappa}}(V)=\sqrt{4\pi V-\kappa V^2}$.
Fill the holes in each connected component, retaining its outer boundary.
The resulting smooth disk has at least the original area and no greater
boundary length. The disk inequality and monotonicity therefore bound
the original component's boundary length below by the model profile at
its original area. Summing and using subadditivity proves the comparison
for the original domain, even when the filled disks overlap.

\subsection{Smooth approximation in a fixed neighborhood}

Recall that $BV(M)$ consists of integrable functions with finite
distributional variation. Strict convergence means $L^1$ convergence
together with convergence of the total masses $|Du|(M)$. For
$u=\chi_E$, that mass is the full ambient perimeter.

The following form of approximation preserves ambient perimeter. No
correction to impose exactly equal volumes is needed.

\begin{lemma}[Smooth approximation]\label{lem:smooth-bv-approximation}
Let $M$ be a complete connected smooth Riemannian manifold without
boundary. Let $E$ have finite ambient perimeter, and suppose that
$|E\setminus K|=0$ for a compact set $K\subset M$. For every open
neighborhood $U\supset K$ with compact closure, there are relatively
compact open sets $E_j\Subset U$ with smooth boundary such that
\[
 |E_j\mathbin{\triangle}E|\longrightarrow0,
 \qquad P(E_j)\longrightarrow P(E).
\]
The sets $E_j$ need not be connected.
\end{lemma}

\begin{proof}
Strict smooth approximation on complete manifolds
\cite[Theorem~2.12 in the author preprint]{GuneysuPallara2015}
provides smooth compactly supported functions $f_j$ with
\[
 \|f_j-\chi_E\|_{L^1(M)}\longrightarrow0,
 \qquad \int_M|\nabla f_j|\,d\vol_g\longrightarrow P(E).
\]
The cited result assumes completeness, smoothness, connectedness, and
absence of boundary; it imposes no curvature condition. We may take real
parts, since doing so cannot increase the $L^1$ error or gradient
integral, and lower semicontinuity gives the reverse limiting bound.

Choose $\eta\in C_c^\infty(U)$ with $0\le\eta\le1$ and $\eta=1$
near $K$, and put $u_j=\eta f_j$. Then $u_j\to\chi_E$ in $L^1(M)$,
and $\chi_E\nabla\eta=0$ almost everywhere. Hence
\[
 \int_M|\nabla u_j|\,d\vol_g
 \le\int_M|\nabla f_j|\,d\vol_g
       +\|\nabla\eta\|_\infty\|f_j-\chi_E\|_1.
\]
Lower semicontinuity shows that the left side also converges to $P(E)$.
Thus all approximating functions now have support in the same compact
subset of $U$.

Write $e_j=\|u_j-\chi_E\|_1$, and choose $0<\delta_j<1/2$ with
$\delta_j\to0$ and $e_j/\delta_j\to0$. For
$\delta_j<t<1-\delta_j$, pointwise comparison with $\chi_E$ gives
\begin{equation}\label{eq:threshold-volume-error}
 \bigl|\{u_j>t\}\mathbin{\triangle}E\bigr|
 \le\delta_j^{-1}e_j.
\end{equation}
The ordinary coarea formula for the smooth function $u_j$ gives
\[
 \int_{\delta_j}^{1-\delta_j}P(\{u_j>t\})\,dt
 \le\int_M|\nabla u_j|\,d\vol_g.
\]
By Sard's Theorem, almost every $t$ is a regular value. We can therefore
choose such a $t_j\in(\delta_j,1-\delta_j)$ with
\[
 P(\{u_j>t_j\})
 \le\frac{1}{1-2\delta_j}\int_M|\nabla u_j|\,d\vol_g+\frac1j.
\]
Set $E_j=\{u_j>t_j\}$. Its boundary is smooth and compact in $U$.
Equation~\eqref{eq:threshold-volume-error} gives $L^1$ convergence of
the characteristic functions. The displayed perimeter bound and lower
semicontinuity give the claimed perimeter convergence.
\end{proof}

Apply Lemma~\ref{lem:smooth-bv-approximation} to a bounded
finite-perimeter set in a two- or three-dimensional Cartan--Hadamard
manifold. The smooth comparison just proved applies to every $E_j$,
including its disconnected cases. Since $|E_j|\to|E|$ and
$\mathcal I_b$ is continuous, it passes to $E$.
Together with Theorem~\ref{thm:normalized-comparison}, this establishes
the normalized comparison for bounded finite-perimeter sets in every
dimension $n\ge2$.

\subsection{Removing boundedness}

We record the cutoff argument separately, since it applies to any
continuous increasing perimeter lower bound.

\begin{lemma}[Passage to finite volume]\label{lem:finite-volume-cutoff}
Let $M$ be a complete connected smooth Riemannian manifold without
boundary. Let $J:[0,\infty)\to[0,\infty)$ be continuous and increasing,
with $J(0)=0$. Suppose that $P(F)\ge J(|F|)$ for every bounded
finite-perimeter set $F\subset M$. Then the same inequality holds for
every measurable $E$ of finite volume and finite ambient perimeter.
\end{lemma}

\begin{proof}
We use two elementary BV identities. For a compactly supported
Lipschitz function $\varphi$ and bounded $u\in BV(M)$, the product rule is
\[
 D(\varphi u)=\varphi\,Du+u\nabla\varphi\,d\vol_g.
\]
It follows from the distributional product rule for smooth multipliers
by smoothing $\varphi$ locally in charts: the approximations converge
uniformly, and their gradients converge in $L^1$ on the fixed compact
support. The boundedness of $u$ permits passage in the second term.
For $0\le u\le1$, the coarea identity is
\begin{equation}\label{eq:bv-coarea}
 |Du|(M)=\int_0^1P(\{u>t\})\,dt.
\end{equation}
The integrand is measurable: the defining perimeter supremum can be
taken over a countable dense family of compactly supported test fields.
For completeness, apply the strict smooth approximation theorem
\cite[Theorem~2.12 in the author preprint]{GuneysuPallara2015} to
$u\in BV(M)$. It applies to arbitrary integrable functions.
The ordinary coarea formula then proves the inequality ``$\ge$'': take a subsequence whose
level sets converge in $L^1$ for almost every $t$, and use lower
semicontinuity and Fatou's Lemma. Such a subsequence exists because
the integral over $t\in\R$ of the $L^1$ distance between the two level
indicators equals the $L^1$ distance between the functions. For the
reverse inequality, write $u=\int_0^1\chi_{\{u>t\}}\,dt$, test against
$\divg X$ for $X\in C_c^\infty(TM)$ with $|X|\le1$, and take the supremum
in the definition of total variation. This also shows that almost every
level in \eqref{eq:bv-coarea} has finite ambient perimeter.

Fix $o\in M$ and, for $R>0$, define
\[
 \varphi_R(x)=\min\{1,\max\{0,R+1-d(o,x)\}\}.
\]
Completeness makes its support compact. It equals one on $B_R(o)$,
vanishes outside $B_{R+1}(o)$, and has Lipschitz constant at most one.
The product rule and \eqref{eq:bv-coarea} give
\begin{equation}\label{eq:cutoff-coarea-perimeter}
 \begin{split}
 \int_0^1P(\{\varphi_R\chi_E>t\})\,dt
 &=|D(\varphi_R\chi_E)|(M)\\
 &\le\int_M\varphi_R\,d|D\chi_E|
       +\int_E|\nabla\varphi_R|\,d\vol_g\\
 &\le P(E)+|E\setminus B_R(o)|.
 \end{split}
\end{equation}
Almost every level is a bounded finite-perimeter set, and for every
$0<t<1$ its volume is at least $v_R=|E\cap B_R(o)|$. The hypothesis
and monotonicity of $J$ therefore bound the left side of
\eqref{eq:cutoff-coarea-perimeter} below by $J(v_R)$.
Since $v_R\to|E|$ and $|E\setminus B_R(o)|\to0$, continuity of $J$
proves the assertion.
\end{proof}

\begin{proof}[Proof of Theorem~\ref{thm:main}]
Null sets have zero perimeter and $\mathcal I_b(0)=0$.
The preceding bounded-set comparison and
Lemma~\ref{lem:finite-volume-cutoff}, with $J=\mathcal I_b$, prove the
assertion for all $n\ge2$ and normalized curvature parameters
$b\in\{0,1\}$.

For an arbitrary negative curvature bound, put $b=\sqrt{-\kappa}>0$
and $\widetilde g=b^2g$. Then
\[
 \Sec_{\widetilde g}=b^{-2}\Sec_g\le-1,
 \qquad |E|_{\widetilde g}=b^n|E|_g,
 \qquad P_{\widetilde g}(E)=b^{n-1}P_g(E).
\]
The last identity holds for ambient BV perimeter: constant metric
rescaling leaves divergence unchanged, while
$d\vol_{\widetilde g}=b^n d\vol_g$ and $|X|_{\widetilde g}=b|X|_g$.
In the dual definition of perimeter, the substitution $Y=bX$ therefore
multiplies the supremum by $b^{n-1}$.
The model functions have the matching relations
\[
 \mathcal V_b(r)=b^{-n}\mathcal V_1(br),\qquad
 \mathcal A_b(r)=b^{-(n-1)}\mathcal A_1(br),\qquad
 \mathcal I_b(V)=b^{-(n-1)}\mathcal I_1(b^nV).
\]
Applying the normalized result to $\widetilde g$ and dividing by
$b^{n-1}$ gives the claimed inequality for $g$ and $\kappa$.
The case $\kappa=0$ was proved directly, without a curvature limit.
\end{proof}
